%% ============================================================
%% Chapter 1 Solutions
%% ============================================================

\section{Chapter 1: Introduction and Overview}

%% ----------------------------------------------------------
\begin{solution}{1.1}{L2}

\solanswer

A \textbf{parameter of interest} (\POI) is an input to the model
whose value may be uncertain or whose influence we wish to quantify.
A \textbf{quantity of interest} (\QOI) is an output of the model
that we care about --- the number we are trying to predict or understand.

\smallskip
For an antibiotic-resistance model in a hospital, two \POIs might be:
the probability that a patient receives an antibiotic in a given
hospital stay, and the probability that a resistant strain is
transmitted per nurse-patient contact. Two \QOIs might be:
the fraction of patients colonized with a resistant strain after
30 days, and the mean length of stay for infected patients.

\smallskip
Running the model tells us the \QOI at the nominal parameter values.
Sensitivity analysis tells us \emph{how much the \QOI would change}
if a \POI were different --- that is, which parameters matter most
for the model's predictions, so we can focus data collection and
intervention design on those parameters.

\solnote{Accept any reasonable \POI/\QOI pair consistent with the
antibiotic-resistance context. The key distinction: \POIs are inputs
controlled or measured by the investigator; \QOIs are outputs that
answer the policy or scientific question. Students sometimes confuse
state variables with \POIs; clarify that parameters govern the
dynamics, while state variables evolve according to those parameters.}

\end{solution}

%% ----------------------------------------------------------
\begin{solution}{1.2}{L2}

\solanswer

\textbf{Forward SA} answers ``What if?'' questions: if this \POI changes,
how does the \QOI respond? It computes sensitivity by varying each
\POI one at a time and observing the effect on the \QOI. The cost
scales with the number of \POIs (one computation per \POI).

\textbf{Adjoint SA} answers ``How come?'' questions: given that the
\QOI has the value it does, which parameters are responsible?
It computes sensitivity by working backward from the \QOI to the \POIs.
The cost scales with the number of \QOIs (one computation per \QOI).

\smallskip
Classifying the four questions:

\solpart{a} \emph{If the mean latent period increases by 10\%, how does
peak infection count change?} --- \textbf{Forward} question (one \POI,
asking about effect on \QOI).

\solpart{b} \emph{The peak infection count exceeded the hospital capacity.
Which parameters are most responsible?} --- \textbf{Adjoint} question
(one \QOI, asking which \POIs drove it).

\solpart{c} \emph{We have 20 uncertain parameters. Which 3 are worth
measuring precisely?} --- \textbf{Forward} (or global) SA: compute SIs
for all 20 parameters and identify the largest.

\solpart{d} \emph{We want $\partial J/\partial p$ for 50 parameters where
$J$ is a single regulatory threshold.} --- \textbf{Adjoint}: one \QOI,
many \POIs, so adjoint is efficient.

\solnote{Part (d) is a good opportunity to preview the cost comparison
(Chapter 5): adjoint gives all 50 sensitivities at the cost of one
computation, vs.\ 50 computations for forward SA.}

\end{solution}

%% ----------------------------------------------------------
\begin{solution}{1.3}{L4}

\solanswer

\solpart{a} \emph{Climate model, 500 parameters, identify 10 most
influential on mean surface temperature, 8 hours per run.}

\textbf{Morris screening} (or adjoint, if available). With 8 hours
per run, a full Sobol analysis at $N = 2000$ would cost $N(m+2) = 1{,}004{,}000$
runs, or about 916 years of computation. Morris screening with
$r = 10$ trajectories costs $r(m+1) = 5010$ runs ($\approx$4 years)
--- still expensive, but feasible on a cluster. Adjoint SA would
cost the equivalent of one forward run plus storage, giving all
500 sensitivities simultaneously. If the adjoint is available,
use it; otherwise, Morris screening on a cluster is the practical alternative.

\solpart{b} \emph{Pollution transport model, 6 parameters, threshold
exceedance at one monitoring station.}

\textbf{Adjoint SA.} The \QOI is a single threshold value at one
location --- exactly the scenario where the adjoint is optimal.
With 6 parameters and 1 \QOI, adjoint costs 1 backward solve vs.\
6 forward perturbations. The adjoint gives the complete gradient
of the threshold value with respect to all 6 parameters simultaneously.

\solpart{c} \emph{Pharmacokinetic model, 8 parameters, patient outcome
is AUC (area under the drug-concentration curve), nonlinear dose-response.}

\textbf{Global SA (Sobol indices).} The dose-response is nonlinear,
so local SA (forward or adjoint) will miss nonlinear contributions
and interactions. Sobol indices capture the full variance-based
sensitivity across the parameter uncertainty range. With 8 parameters,
a Sobol analysis at $N = 2000$ requires $2000 \times 10 = 20{,}000$ evaluations,
which is typically feasible for a pharmacokinetic ODE model.

\solpart{d} \emph{Structural engineering FEM model, 200 parameters,
single safety factor \QOI, model takes 10 seconds per run.}

\textbf{Adjoint SA.} Single \QOI (safety factor), 200 parameters,
and 10 seconds per run: Sobol at $N=2000$ would take $2000\times202 = 404{,}000$
evaluations $\approx 47$ days. Adjoint gives all 200 sensitivities
in $\approx$20 seconds (one forward + one backward solve). The
adjoint is the clear choice when the model is differentiable and
$m \gg r$.

\solnote{This exercise builds intuition for the decision table
(Table 4.2 in the text). Common student errors: choosing global SA
for (b) because ``nonlinearity might be present'' (the model is not
stated to be nonlinear, and the threshold \QOI is a single value
at one station, so adjoint is correct); or choosing forward SA for
(d) without computing the cost.}

\end{solution}

%% ----------------------------------------------------------
\begin{solution}{1.4}{L5}

\solanswer

\textbf{Model answer} (acceptable answers will vary in specifics but
should include all four required components):

\medskip
\noindent\textit{SA plan for a dengue fever model with 9 parameters.}

\textbf{QOI choice.} I would use peak infectious count $\max_t I(t)$ as the
primary \QOI, because it determines whether the healthcare system is
overwhelmed --- the operationally critical threshold. A secondary \QOI
would be the basic reproduction number $R_0$ (analytically tractable
and directly related to whether an epidemic occurs). The relationship
between the two provides a consistency check.

\textbf{Identifying the 3 parameters to measure.}
Run a local SA using the forward sensitivity equations (analytic
if the model is an ODE, or centered finite differences otherwise)
to compute all 9 sensitivity indices $\SI{p_j}{\text{peak}~I}$.
Rank by $|\SI{p_j}{\cdot}|$. The three parameters with the largest
absolute SI values are the priority for precise measurement.
If two parameters appear in the model only as a product (structural
correlation), treat the product as a single effective parameter.

Before reporting, also run a Morris screening (15 trajectories,
cost: $15 \times 10 = 150$ model evaluations) to flag any parameters
where nonlinearity makes the local SI misleading.

\textbf{Additional information needed for intervention design.}
SA tells us which parameters \emph{matter most for predictions},
not which can be controlled by public health interventions.
I would cross-reference the top-ranked parameters with a table of
interventions that affect each parameter (e.g., insecticide spraying
affects vector contact rate; surveillance programs affect reporting rate).
Only parameters that are both high-sensitivity \emph{and} actionable
with available resources should drive intervention recommendations.

\textbf{Limitation.}
The most important limitation is that local SA is a linearization:
the sensitivity indices are derivatives at the nominal parameter values.
If the model is highly nonlinear (e.g., near a threshold between
epidemic and no-epidemic regimes), local SA may not accurately
reflect sensitivity across the full plausible parameter range.
I would flag this to the team and recommend a global SA (Sobol)
on the top 4--5 parameters identified by local SA, to check for
nonlinear effects.

\solnote{Grading rubric: 1 point for \QOI justification; 2 points
for the parameter-identification procedure (must mention a specific
SA method and be operationally specific about the ranking); 1 point
for the intervention-actionability distinction; 1 point for identifying
a genuine limitation of local SA. Deduct 1 point if student recommends
global SA for all 9 parameters without noting the cost implication
(dengue models are often ODEs; full Sobol is feasible, so this is
not wrong, but the question asks them to use SA to guide \emph{field}
measurement decisions, which requires identifying the top 3 specifically).}

\end{solution}

%% ----------------------------------------------------------
\begin{solution}{1.5}{L6}

\solanswer

\textbf{Problem 1: Conflation of ``all parameters simultaneously''
with global SA.}

Computing local sensitivity indices (derivatives at best-fit values)
does \emph{not} analyze all parameters simultaneously. Local SA
perturbs one parameter at a time, holding all others fixed.
The statement implies that the joint effect of simultaneous parameter
uncertainty was assessed, but it was not.
What should have been done: either (a) state that local SA was used,
which analyzes one parameter at a time and captures only linear,
individual effects; or (b) if the goal was to assess robustness to
joint uncertainty, perform global SA (e.g., Sobol indices), which
does evaluate sensitivity across the full multi-dimensional parameter
space.

\textbf{Problem 2: ``Robust to parameter uncertainty'' is not a
valid conclusion from local SI values.}

Sensitivity indices measure the \emph{rate of change} of the \QOI
with respect to each parameter. A model with all small SIs at the
best-fit values is locally insensitive near that point, but could
be highly sensitive elsewhere (e.g., near a bifurcation).
``Robust'' should mean that the \QOI changes little across the full
plausible range of parameter values --- a global property, not a
local one. What should have been done: specify a plausible parameter
uncertainty range for each parameter, compute the resulting range
of \QOI values, and only then make a robustness claim.

\textbf{Optional additional problem} (award bonus credit): Local SA
at ``best-fit parameter values'' is circular if the model was fitted
to the same data being used to validate it. Sensitivities may be
artificially small near the best-fit point if the fitting procedure
has forced the model to be locally insensitive to parameter changes.

\solnote{This is a Bloom L6 synthesis exercise. Full credit requires
identifying problems that go beyond surface-level criticism. The
strongest answers connect ``local vs.\ global'' to the mathematical
definition of the SI (derivative vs.\ variance decomposition) rather
than simply noting vague imprecision. Partial credit (1 point each)
for correct identification of any two genuine methodological issues
with specific alternative recommendations.}

\end{solution}
